\chapter{Calibrazione e limiti di simulazione}

Questa sezione descrive aspetti di calibrazione software e limiti della simulazione.

\section{Calibrazione software}
\begin{itemize}
  \item \textbf{Temperatura}: il calcolo NTC usa parametri standard (BETA, RL10). La precisione reale dipende dal sensore e dal cablaggio, ma in simulazione il modello e stabile.
  \item \textbf{Umidita}: il DHT22 in simulazione restituisce valori ideali; in ambiente reale e consigliata una calibrazione empirica.
  \item \textbf{Lux}: la conversione da fotoresistenza e approssimata; il target in lux e un riferimento software piu che una misura certificata.
\end{itemize}

\section{Limiti della simulazione}
\begin{itemize}
  \item Il PWM e semplificato (valore pieno quando serve luce).
  \item Il tempo di risposta dei sensori e idealizzato.
  \item Alcuni fenomeni reali (rumore, deriva, latenza) non sono riprodotti.
\end{itemize}
